\SourceLeaf{264}%
\[
\frac{d\,dy}{dx^{2}} = -\alpha\alpha\,\frac{1-\dfrac{dy^{2}}{\alpha\alpha\,dx^{2}}}{y},
\]

\[
\frac{y\,ddy}{dx^{2}} - \frac{dy^{2}}{dx^{2}} + \alpha\alpha = 0.
\]

\noindent Multiplicirt mit $\dfrac{2\,dy}{y^{3}}$:

\[
\frac{2\dfrac{dy}{dx}\cdot d\dfrac{dy}{dx}}{yy} - \frac{2\left(\dfrac{dy}{dx}\right)^{2}dy}{y^{3}} + \frac{2\alpha\alpha\,dy}{y^{3}} = 0.
\]

\noindent Integrirt

\[
\frac{1}{yy}\left(\frac{dy}{dx}\right)^{2} - \frac{\alpha\alpha}{yy} = \text{Const.} = -kk,
\]
\[
-kkyy + \alpha\alpha - \left(\frac{dy}{dx}\right)^{2} = 0,
\]
\[
dx = \frac{dy}{\BareRoot{(\alpha\alpha-kkyy)}} = \frac{\dfrac{k}{\alpha}dy}{k\BareRoot{\left(1-\dfrac{kk}{\alpha\alpha}yy\right)}},
\]
\[
kx = \text{Arc.}\sin\frac{k}{\alpha}y + \text{Const.},
\]
\[
y = \frac{\alpha}{k}\sin(kx+\lambda),
\]
\[
fx = \frac{\alpha}{k}\sin kx,
\]
\[
gx = \cos kx.
\]

\smallskip
\begin{center}

\end{center}

\begin{center}
\textsc{Bemerkungen.}
\end{center}

\noindent \textsc{Gauss} hat die vorstehende Notiz auf einem Zettel (4 Seiten 8\textsuperscript{0}) niedergeschrieben, der sich als Einlage in seinem Exemplare der \emph{»Geometrischen Untersuchungen zur Theorie der Parallellinien} von Nicolaus \textsc{Lobatschewsky}, Berlin 1840« vorgefunden hat; als Abfassungszeit darf daher 1840 bis 1846 angenommen werden.

Der Zweck der Notiz ist, unter der Voraussetzung, dass f\"ur Dreiecke, deren drei Seiten unendlich klein sind, die Euklidische Geometrie gilt, die Gleichungen herzuleiten, die zwischen den Seiten und Winkeln eines endlichen Dreiecks bestehen; es gen\"ugt, diese Beziehungen f\"ur ein rechtwinkliges Dreieck herzuleiten. Die Durchf\"uhrung dieses Gedankens gestaltet sich auf Grund des vorstehenden Textes folgendermassen.
